% Source: https://artofproblemsolving.com/community/c3h2319369_100th_post_problems
% Stable identity: aops-topic-2319369-q75
The orbit has radius $6400+2400=8800$ km and circumference $2\pi\cdot8800=17600\pi$ km. In $1$ hour the satellite covers a third of it:
\[\dfrac{17600\pi}{3}\approx\boxed{18400\text{ km}}\ (18431\text{ km}).\]
It has turned through $120^\circ$, so its straight-line distance from the start is the chord $2\cdot8800\sin60^\circ=8800\sqrt3\approx\boxed{15200\text{ km}}$ ($15242$ km).

\begin{asy}
usepackage("amsmath");
size(5cm);
defaultpen(fontsize(10pt));
pair O=(0,0);
real R=6.4, r=8.8;
filldraw(circle(O,R),lightgray,black);
draw(circle(O,r),blue);
pair P=r*dir(90), Q=r*dir(210);
draw(arc(O,r,90,210),red+linewidth(1.4));
draw(O--P,dashed); draw(O--Q,dashed); draw(P--Q,dotted);
draw(arc(O,1.5,90,210));
label("$120^\circ$",(-2,2),fontsize(8pt));
dot(P); dot(Q); dot(O);
label("start",P,N); label("after $1$ h",Q,SW);
label("Earth",(0,-3.2)); label("$6400$",(0,3.6),E,fontsize(8pt)); label("$2400$",(0,7.6),E,fontsize(8pt));
\end{asy}
